\documentclass[a4paper, twoside]{ltjsarticle}

%%%\usepackage[top=13truemm,bottom=8truemm,]{geometry}
%\setlength{\headheight}{10truemm}
%\setlength{\topmargin}{-25truemm}%% \topmarignを-0.5インチにする
%\setlength{\headsep}{0pt}
%\setlength{\marginparsep}{0pt} %本文領域と傍注領域との空き
%\setlength{\footskip}{0truemm}
%%\usepackage[hiragino-pro]{luatexja-preset}
\usepackage[hiragino-pron,jis2004]{luatexja-preset}
%%\usepackage[dvipdfmx]{graphicx}
\usepackage{graphicx}
\usepackage{wrapfig}
\usepackage{url}

\title{「ゴシック･アンド･ロリータ」東京公演 } 
\author{ 出演 成田達輝 (ヴァイオリン), 山根明季子 (作曲), 梅本佑利(作曲) \thanks{mumyo} }
\date{2023年4月2日}

\begin{document}

\maketitle

\begin{abstract}
ゴシック・アンド・ロリータ（ゴスロリ）は、ヨーロッパの伝統を下地に日本のストリートで再構築された服飾文化、サブカルチャー。タナトスを想起させるゴシックと、少女の典型であるロリータが独自に結び付いた表現である。本コンサートでは、ヨーロッパの伝統音楽、ヴァイオリン独奏曲のフォーマットを用いて、ゴスロリと芸術音楽を接続する。
\end{abstract} 

\section*{公演の概要}
\begin{itemize}
\item 2023年4月2日 (日) 開催 
\item 会場: BUoY（東京・北千住）
\item 主催: mumyo / 合同会社無名
\item 第1回：開演 14:00（開場 13:30）\\
 トークゲスト：朝藤りむ（ファッションブランド pays des fées デザイナー）
\item 第2回：開演 17:00（開場 16:30）\\
 トークゲスト：NABEchan（イラストレーター）　
\item 一般4000円、25歳以下1500円
\end{itemize}

現代音楽の分野で少女性・日本のポップ・サブカルチャーなどのテーマに取り組んできた作曲家の山根明季子、そして梅本佑利の全新作8作品によるプログラム。国内外主要オーケストラと多数共演しているヴァイオリニスト・成田達輝のストラディヴァリウスによる演奏でお届けする。

各公演の前または後にトークを実施。14時公演ではファッションブランド pays des fées デザイナーの朝藤りむ、17時公演ではイラストレーターのNABEchan(本公演メインビジュアルを担当)をゲストに迎える\footnote{ 各公演とも、ゲスト以外のプログラムは同一です。}。

\subsection*{曲目}

\begin{description}
\item[山根明季子]「黒いリボンをつけたブーレ」,「リボン集積」,「リボンの血肉と蒸気」,「パニエ、美学」（2022）
\item[梅本佑利]「Embellish Me!」,「Melt Me!」,「コピー・アンド・ペースト，大量生産／消費された不規則／不完全な形状のプラスチック真珠そして私。」,「廃墟・秋葉原のアリス 1, 2」（2022）
\end{description}

\samepage
\rightline{ \gt (裏面に続く)}
\newpage

\section*{ステートメント全文　「少女、アリス、山根明季子」}

山根明季子にとっての最初の委嘱作品である2008年作曲『ケミカルロリイタ』（チューバとピアノのための）\footnote{ 橋本晋哉による委嘱。}では、チューバがルイス・キャロル『不思議の国のアリス』の一節を語る。今回の「ゴシック・アンド・ロリータ」公演で演奏される、梅本による『～Me』（2022, ヴァイオリンのための）の連作において、不思議の国のアリスが引用されるのは、山根の出発点である『ケミカルロリイタ』の系譜を明確にするためでもある。その後の作品においても、山根は、『少女メランコリー』（2011, ヴァイオリンとトイピアノのための）、『ハラキリ乙女』（2012, 琵琶とオーケストラのための）などで、積極的に「少女」を描いてきた。

そんな山根の描く「少女」はどこか憂鬱だ。彼女の作品には、少女の自傷的な衝動が見て取れる。『少女メランコリー』では、「破壊的でグロテスクな思考回路」\footnote{\url{ https://akikoyamane.com/post/169459501812/girl-melancholy}}が描かれ、『ハラキリ乙女』では、「カッターなどの鈍い刃物」\footnote{\url{https://akikoyamane.com/post/168038351977/harakiri-maiden}}で、自ら肌を切り付ける。最初の委嘱作品の表記にみられるように、その少女は「ロリータ」ではなく、「ロリイタ」と、日本のロリータファッションの文脈であることが伺える。漂う死の匂いとロリイタ。そんな彼女の作品は、いかにもゴスロリ的なのである。

日本のゼロ年代思想や80年代以降のオタク論のように、言説の世界において、サブカルチャーとしての「ゴスロリ」は、「オタク」ほどには語られていない。山根は、今回のプロジェクトに関わるまで --- 自身がゴスロリ的な作品を作曲していたとはほとんど認識しておらず、自己を批評的に言語化してこなかった --- と梅本に語っているが、その内向的な性格は、まさに「ゴスロリの精神」のようである。オタク（の一部）が、現代美術や思想、言説をまとって、ある種の社交性を身に付けたのに比較して、ゴスロリは、メディア、言論への登壇を極力避けた。言葉を発しないということは、極めて強力な「防御」として、「美」を死守する手段なのかもしれない。

西洋音楽における日本式ポップアートを開拓した最重要人物である山根明季子が、長年の的外れな評論で、ほとんどまともな言語化をされることなくここまで来てしまったのには、上述のような一筋縄ではいかない複雑な要因があるのだろう。だがそれと同時に、私は絶対に、この価値ある革新の文脈を埋もれさせたくはない。なぜならば、私の、そして未来の芸術音楽の源流が間違いなくそこにあるからだ。そして、そのために、いま言語化が必要なのだ。我々 (mumyo) の活動の第一段階に「ゴシック・アンド・ロリータ」を置いたのは、まずそのエニグマを芸術でもって解き明かし、提示することにある。

ゴスロリを通して過去を顧みて、未来の世界を占う。情報の樹海を彷徨う現代の少女は、あらゆる装飾で防御し、その美でもって死を超越する。生まれもった肉体とあらゆる物質を自在に縫い合わせ、無限の可能性を持って変容すること。生まれ持ったものが物事の本質であるなどと、本質主義的に捉える時代は終わった。出自について語る意味と、出自の無意味さは共存する。書けば書くほど、聴けば聴くほどナンセンス。インベリッシュ・ミー！　アリスの体は知でもって、その過剰な装飾によって変身する。 \qquad (梅本佑利)

\subsubsection*{2023年4月2日 (日) 開催「ゴシック･アンド･ロリータ」}
%%\begin{figure}[hbtp]
 %%\centering
%% \includegraphics[keepaspectratio, width=3cm]{gl_qrcode.eps}
 %%\caption{チケットは、イープラスにて発売中}
 %%\label{fig.1}
%%\end{figure}

\begin{wrapfigure}{r}[0pt]{0.2\textwidth}
 \centering
   \includegraphics[width=3cm,clip]{gl_qrcode.png}
 %\caption*{チケット販売}
\label*{チケット販売}
\end{wrapfigure}

チケットは、イープラスにて発売中\\
{\tt https://eplus.jp/sf/detail/3825110001-P0030001}


\end{document}
